\originalchapter{8.421公开作业六至十}
本章题源仍为Spring 2012原件, 解答为编者补写. 辐射模式简并、饱和参数与线宽约定在各题单独说明.
\originalsection{作业六: 黑体与腔}
\source{Assignment 6, 第1题.}
\begin{exercise}(a) Na主跃迁590 nm, 寿命16 ns. 假定每次黑体吸收损失一个磁阱原子, 要保持至少一分钟, 求每模光子数和相应温度, 是否需屏蔽300 K或冷到4 K? (b) 氢$1S,F=1\to F=0$的21 cm线由什么多极耦合, 用题给电偶极$ea_0$, 磁偶极$\mu_B$, 电四极$ea_0^2$估寿命. (c) 求21 cm在300 K、4 K的每模光子数、诱导向下跃迁率与一分钟实验的影响; (d) 为什么更多每模光子并不意味着本例中更短寿命?\end{exercise}
\begin{solution}(a) 在等权二态简化下$R_{\rm abs}=\Gamma\bar n$, 所以$\bar n\le16\times10^{-9}/60=2.67\times10^{-10}$. $T=hc/[k_B\lambda\ln(1+1/\bar n)]\simeq1.11\times10^3$ K. 300 K的光学黑体占据远小于此值, 无需为这条跃迁冷却真空腔; 若精确计激发多重态, 限值会有简并因子, 结论不变. (b) 两个$1S$态同宇称, 电偶极禁戒, 主导为磁偶极. $\Gamma_M=\mu_0\omega^3|\mu|^2/(3\pi\hbar c^3)$, 取$\mu_B$得约$2.9\times10^{-15}$ s$^{-1}$, 寿命约$3.5\times10^{14}$ s, 约千万年. 题给电四极矩即便形式存在也多出$(ka_0)^2$抑制, 且$F=1\to0$不满足二级张量三角条件. (c) $hc/(k_B\lambda)=0.0685$ K, 故$\bar n_{300}\simeq4378,\bar n_4\simeq57.9$. 受激向下率$\Gamma_M\bar n$约$1.3\times10^{-11},1.7\times10^{-13}$ s$^{-1}$, 一分钟影响可忽略. (d) 发射率还乘模式密度和耦合矩阵元, 磁偶极弱且频率立方极小; 光子占据只是一个因子.
\end{solution}
\source{Assignment 6, 第2题.}
\begin{exercise}二能级上态原子放入无损单模空腔, 腔频$\omega_c\ne\omega_0$. (a) 是否有定态? (b) 求仍处上态的$P_e(t)$; (c) 总能量是否守恒?\end{exercise}
\begin{solution}(a) $\ket{e,0},\ket{g,1}$的二乘二Hamiltonian有两个缀饰本征态, 所以有定态, 但单独$\ket{e,0}$不是. (b) 记$\delta=\omega_c-\omega_0$, $Q=\sqrt{\delta^2+4g^2}$, $P_e=1-[4g^2/Q^2]\sin^2(Qt/2)$. (c) 时间独立完整Hamiltonian的期望值严格守恒. 原子和光子的裸能量若失谐并不相等, 其变化由交互作用能补偿; 不可只把裸原子能量与裸腔能量相加而忽略耦合项.
\end{solution}
\source{Assignment 6, 第3题.}
\begin{exercise}Rydberg相邻$n\to n-1$跃迁角频率原子单位约$n^{-3}$, 电偶极约$n^2$. (a) 求$n=50\to49$频率Hz和可调到该频率的最小闭合立方腔尺寸; (b) 求真空Rabi频率Hz.\end{exercise}
\begin{solution}(a) 以题给量级$\omega=(E_h/\hbar)n^{-3}$, $f=(E_h/h)n^{-3}\simeq52.6$ GHz. 精确氢能差给54.3 GHz, 是近似尺度的误差范围. 理想PEC立方腔最低电磁模有两个非零指数, $f=c/(\sqrt2L)$, 故$L\simeq4.03$ mm. 不能用允许三指数全零的虚假腔模. (b) 选原子在电场反节点, 偶极沿该模电场, 正弦积模式有$V_{\rm eff}=L^3/4$. $d=2500ea_0$, $g=d\sqrt{\hbar\omega/(2\epsilon_0V_{\rm eff})}/\hbar$, 因而真空布居振荡的角频率$2g$, 普通频率$2g/(2\pi)$约0.703 MHz. 若用均匀场体积$L^3$作量级估计, 值减半. 原子位置与模式归一化必须给定, 题中尺度律不能支持唯一精确实验耦合值.
\end{solution}

\originalsection{作业七: 饱和与弛豫}
\source{Assignment 7, 第1题.}
\begin{exercise}(a) 以振子强度$f$, $\alpha$与$\omega$表达Einstein $A$, 取Na590 nm的单有效Cartesian振子$f=1$估寿命; (b) 把Na当闭合二级系统, 求饱和强度. 题中饱和定义为弱场共振激发率等于自然率的一半.\end{exercise}
\begin{solution}按第4章方向振子强度定义$f=2m\omega d^2/(\hbar e^2)$. 式\eqref{eq:spont}给$A=2\alpha\hbar\omega^2f/(3mc^2)$, Na约$6.39\times10^7$ s$^{-1}$, 寿命约15.7 ns. 闭合二态饱和强度$I_s=\pi hcA/(3\lambda^3)$约64.7 W/m$^2$, 即6.47 mW/cm$^2$. 在$s=1$的真实稳态散射率为$A/4$, 题中$A/2$指未计饱和的激发率, 不应混同.
\end{solution}
\source{Assignment 7, 第2题.}
\begin{exercise}宽带光驱动等权两态, 受激吸收与发射率$R$, 自发率$\Gamma$, $s=2R/\Gamma$. (a) 求平衡激发/基态布居比; (b) 求自发率及漂白截面; (c,d) 求$s=1$的每角频率能量密度并说明偶极抵消、每模占据; (e) 共振Lorentz激光总强度$I_0$, 全宽$\Gamma'\gg\Gamma$, 求中心能量密度及$I_s$; (f) 同强度单色场Rabi频率$\omega_R$, 证$R=\omega_R^2/\Gamma'$并给饱和值; (g) 为什么$\Gamma'=\Gamma$延伸给单色饱和表达式?\end{exercise}
\begin{solution}(a) $\dot p_e=R(1-p_e)-(R+\Gamma)p_e$, 故$p_e/p_g=s/(s+2)$. (b) $p_e=s/[2(1+s)]$, $R_{\rm sc}=\Gamma s/[2(1+s)]$, 除以光子流得$\sigma(s)=\sigma(0)/(1+s)$. (c,d) $R=Bu_\omega$, $u_{\omega,s}=\Gamma/(2B)=\hbar\omega^3/(2\pi^2c^3)$, 因$A/B$无偶极, 所以$\bar n=1/2$. (e) 归一化Lorentz线中心为$2/(\pi\Gamma')$, $u_\omega(\omega_0)=2I_0/(\pi c\Gamma')$. 解得$I_s=\hbar\omega^3\Gamma'/(4\pi c^2)$, 偏振、传播方向和简并需与定义$B$时一致; 若指定单一偏振闭合跃迁, 其有效$B$与各向同性热场系数不同. (f) 在相同偏振约定下, 相关时间$2/\Gamma'$的场黄金律给$R=\omega_R^2/\Gamma'$, 所以$s=1$时$\omega_R^2=\Gamma\Gamma'/2$. (g) 单色场的线宽为零, 但原子相干寿命使有效频谱重叠宽为$\Gamma$. 因而$\omega_R^2=\Gamma^2/2$与Bloch结果一致; 这是在响应函数中以原子自然宽度取代光宽, 不是把宽带近似$\Gamma'\gg\Gamma$当作在等号处仍严格成立.
\end{solution}
\source{Assignment 7, 第3题.}
\begin{exercise}驱动$\omega_R$, 失谐$\delta$, 布居弛豫$\Gamma_1$, 相干弛豫$\Gamma_2$. (a) 比较远失谐时无阻尼Rabi概率的时均与弱场缀饰态概率; (b) 变换$\widehat\rho_{ab}=\rho_{ab}\e^{-\ii\omega t}$写四个Bloch方程, 解稳态; (c) 无驱动先解密度矩阵, 再从原子--真空--一光子的纯态取环境迹, 证明纯自发辐射$\Gamma_1=\Gamma,\Gamma_2=\Gamma/2$; (d) 写真实稳态及饱和参数, 与(a)猜测比较.\end{exercise}
\begin{solution}(a) 无耗散概率$\omega_R^2\sin^2(\sqrt{\delta^2+\omega_R^2}t/2)/(\delta^2+\omega_R^2)$, 远失谐时均约$\omega_R^2/(2\delta^2)$. 一阶弱场幅度为$\omega_R/(2\delta)$, 缀饰基态激发布居约$\omega_R^2/(4\delta^2)$, 相差2, 说明简单抹掉振荡不是耗散稳态推导.

(b) 四式为$\dot\rho_{aa}=\ii\omega_R(\widehat\rho_{ab}-\widehat\rho_{ba})/2+\Gamma_1\rho_{bb}$,$\dot\rho_{bb}=-\dot\rho_{aa}$, $\dot{\widehat\rho}_{ab}=(-\ii\delta-\Gamma_2)\widehat\rho_{ab}+\ii\omega_R(\rho_{aa}-\rho_{bb})/2$, $\dot{\widehat\rho}_{ba}$为其共轭. 先解相干再代入布居, 得$\rho_{bb}=(\omega_R^2\Gamma_2/2\Gamma_1)/[\delta^2+\Gamma_2^2+(\Gamma_2/\Gamma_1)\omega_R^2]$.

(c) 无驱动$\rho_{bb}=\rho_{bb}(0)\e^{-\Gamma_1t}$,$\rho_{aa}=1-\rho_{bb}$,$\rho_{ab}=\rho_{ab}(0)\e^{-\Gamma_2t}$. 联合态$\alpha_a\ket{a,0}+\alpha_b(t)\ket{b,0}+\sum_kc_k\ket{a,1_k}$的原子密度对角为$|\alpha_a|^2+\sum|c_k|^2,|\alpha_b|^2$, 非对角为$\alpha_a\alpha_b^*$. 激发幅度按$\e^{-\Gamma t/2}$下降, 其平方按$\e^{-\Gamma t}$, 所以相干率是一半. (d) $\rho_{bb}=\omega_R^2/[4\delta^2+\Gamma^2+2\omega_R^2]=(s/2)/[1+s+4\delta^2/\Gamma^2]$. 远失谐为$\omega_R^2/(4\delta^2)$, 与弱场缀饰态一致, 与Rabi时均猜测不一致.
\end{solution}

\originalsection{作业八: 线形与偏振}
\source{Assignment 8, 第1题.}
\begin{exercise}两独立展宽线形卷积. (a) Lorentz全宽$\Gamma_1,\Gamma_2$的卷积; (b) Gaussian标准差$\Gamma_1,\Gamma_2$的卷积; (c) 标准差$\Gamma/2$的Gaussian全宽, 全宽$\Gamma$的Lorentz均方根偏差.\end{exercise}
\begin{solution}归一化Lorentz傅里叶变换为$\e^{-\Gamma_i|t|/2}$, 相乘对应全宽$\Gamma_1+\Gamma_2$. Gaussian变换为$\e^{-\Gamma_i^2t^2/2}$, 新标准差为$\sqrt{\Gamma_1^2+\Gamma_2^2}$. Gaussian全宽$\Gamma\sqrt{2\ln2}$. Lorentz尾部为$1/\omega^2$, 二阶矩积分发散, 因而没有有限rms偏差.
\end{solution}
\source{Assignment 8, 第2题.}
\begin{exercise}总自发率$A$, $J=2\to J'=1$. 记$a:2\to1,b:1\to1,c:0\to1,d:1\to0,e:0\to0$. 由每上态总率$A$、每下态收到总率$5A/3$、不偏振总体三偏振等强及$m\to-m$对称性求五率并画分支; 再由Wigner--Eckart或CG系数验证.\end{exercise}
\begin{solution}伸展态只有一支,$a=A$. $b+d=A$,$2c+e=A$, 下态$m'=1$给$a+b+c=5A/3$, $m'=0$给$2d+e=5A/3$. 总$\pi$偏振强度为$2b+e=5A/3$. 联立得$b=d=A/2,c=A/6,e=2A/3$. 以两角动量1相加,$\ket{2,1}=(\ket{1,0}+\ket{0,1})/\sqrt2$, $\ket{2,0}=(\ket{1,-1}+2\ket{0,0}+\ket{-1,1})/\sqrt6$, 其系数平方正给$1/2,1/6,2/3$. 镜像支与对应率相同.
\begin{center}\begin{tikzpicture}[>=Stealth]
\foreach \m in {-2,-1,0,1,2}{\draw (2*\m-.3,2.3)--++(.6,0);\node[above] at(2*\m,2.3){$\m$};}
\foreach \m in {-1,0,1}{\draw (2*\m-.3,0)--++(.6,0);\node[below] at(2*\m,0){$\m$};}
\foreach \u/\v in {-2/-1,-1/-1,-1/0,0/-1,0/0,0/1,1/0,1/1,2/1}{\draw[->] (2*\u,2.23)--(2*\v,.07);}
\node[left] at(-4.4,2.3){$J=2$};\node[left] at(-4.4,0){$J'=1$};
\end{tikzpicture}\end{center}
图中只画允许分支, 避免把速率标签挤在交叉箭头上. 右半边$2\to1$为$A$, $1\to1,0$各$A/2$, $0\to\pm1$各$A/6$, $0\to0$为$2A/3$; 左半边依镜像相同.
\end{solution}
\source{Assignment 8, 第3题.}
\begin{exercise}密度$n$, 共振$\omega_0$, 自然宽$\Gamma$, 光强$sI_s$, 盒长$L$, 薄样品. (a) 静止原子的两布居、吸收截面与吸收比例; (b) 质量$m$, 温度$T$, 沿$z$的泵浦在基态速度分布烧Bennett孔, 求位置、宽度、深度; (c) 扫泵浦频率的Doppler谱及饱和影响; (d) 加弱探针$(k_p,\omega_p)$写吸收积分、窄结构位置与宽度; 再把泵浦反射回来, 比较非零及零失谐的速度孔, 求中央结构宽度. 不要求繁杂积分, 图需体现特征.\end{exercise}
\begin{solution}令$D(v)=1+[2(\delta-k_Lv)/\Gamma]^2$, 热分布$nf(v)=n\sqrt{m/(2\pi k_BT)}\e^{-mv^2/(2k_BT)}$. (a) $n_e=ns/[2(D+s)]$, $n_g=n-n_e$, 净吸收截面$\sigma_0/(D+s)$, 比例$nL\sigma_0/(D+s)$. (b) $n_g(v)=nf(v)[1-s/(2(D(v)+s))]$, 孔心$v=\delta/k_L$, 孔的速度全宽$\Gamma\sqrt{1+s}/k_L$, 中心相对深度$s/[2(1+s)]$.

(c) 吸收为$L\sigma_0\int nf(v)/[D(v)+s]\dd v$. 当Doppler宽远大于$\Gamma\sqrt{1+s}$, 外包络宽度由热分布决定, 饱和主要减小幅度; 极强功率仍可使均匀宽度不可忽略. (d) 弱探针净吸收近似
\[\mathcal A_p=L\sigma_0\int nf(v)\frac{D(v)}{D(v)+s}\frac{\dd v}{1+[2(\delta_p-k_pv)/\Gamma]^2}.\]
这用的是布居差$\rho_g-\rho_e$, 不是仅用$\rho_g$. 窄凹陷在$\delta_p=(k_p/k_L)\delta$, 对反向等波矢是$\delta_p=-\delta$. 把泵浦孔Lorentz半宽$\Gamma\sqrt{1+s}/(2k_L)$与探针响应半宽$\Gamma/(2|k_p|)$卷积, 探针频谱窄凹陷全宽约$\Gamma[1+|k_p/k_L|\sqrt{1+s}]$. 两束等强反射泵浦共同饱和时,$s_{
m tot}(v)=s/[1+4(\delta-kv)^2/\Gamma^2]+s/[1+4(\delta+kv)^2/\Gamma^2]$, $n_g=nf(1+s_{
m tot}/2)/(1+s_{
m tot})$. 非零失谐有$\pm\delta/k$两孔, 零失谐两孔重合, 中心深度$2s/[2(1+2s)]$. 在弱反射探针约定下扫同频$\delta$时, 两孔中心以$2\delta/k$分离, 中央凹陷全宽约$\Gamma(1+\sqrt{1+s})/2$, 弱泵极限$\Gamma$. 若反射光也强, 需用上述共同饱和积分, 不能声称同一简单宽度在任意强度精确成立.
\begin{center}\begin{tikzpicture}\begin{axis}[width=.75\linewidth,height=4cm,xlabel={失谐(示意单位)},ylabel={弱探针吸收},domain=-3:3,samples=150]
\addplot[dashed]{exp(-x*x/4)};\addplot[thick,structurecolor]{exp(-x*x/4)*(1-.5/(1+25*x*x))};
\end{axis}\end{tikzpicture}\end{center}
图只表示宽Doppler包络与窄饱和凹陷, 不以该示意函数拟合上式.
\end{solution}

\originalsection{作业九: 随机场与Bragg转移}
\source{Assignment 9, 第1题.}
\begin{exercise}扰动矩阵元$x(t)\e^{-\ii\omega t}/2$, $x=\pm A$以每方向速率$\Gamma$随机翻转. 原题两态衰减共同给相关振幅因子$\e^{-\gamma|\tau|/2}$. (a) 用初始$x(0)=A$的条件概率$p_\pm$写相关函数; (b) 写速率方程并求$p_+-p_-$; (c) 求相关函数与共振线形.\end{exercise}
\begin{solution}(a) $\langle x(0)x(\tau)\rangle=A^2(p_+-p_-)$. (b) $\dot p_+=-\Gamma p_++\Gamma p_-$,$\dot p_-=-\dot p_+$,$p_\pm=(1\pm\e^{-2\Gamma\tau})/2$. (c) 去掉载频后$G(\tau)=A^2\e^{-(2\Gamma+\gamma/2)|\tau|}/4$. 频谱为$A^2\lambda/[2(\Delta^2+\lambda^2)]$,$\lambda=2\Gamma+\gamma/2$, 全宽$4\Gamma+\gamma$. 这里$\gamma/2$依原题给定衰减约定, 若另把$\gamma$定义为单能级布居损失率, 需先重新确定两态相干衰减, 不能只替换符号.
\end{solution}
\source{Assignment 9, 第2题.}
\begin{exercise}初态$\ket{g,0}$, 中间$\ket{e,\hbar k}$, 末态$\ket{g,2\hbar k}$. 两同偏振反向光$E_1=E_0\cos(kz-\omega t)$,$E_2=E_0\cos[-kz-(\omega+\Delta\omega)t]$, $\Delta\omega<0$, 中间单光子失谐$\Delta$足够大. (a) 求二光子共振条件、带空间因子的交互Hamiltonian、长度$L$周期盒的动量波函数、二光子Rabi频率和有效$i,f$矩阵元; (b) 求光周期平均AC势$U(z,t)$, 算$i,f$矩阵元, 与(a)比较. 内部偶极矩阵元$D_{eg}$.\end{exercise}
\begin{solution}(a) 末态反冲能$4E_R$, 所以$-\Delta\omega=4E_R/\hbar=2\hbar k^2/m$. $H'=-D_{eg}(\ket e\bra g+\hc)(E_1+E_2)$, 两束各自带$\e^{\pm\ii kz}$. 盒中$\phi_p=\e^{\ii pz/\hbar}/\sqrt L$, $p=2\pi\hbar n/L$. 设$\Omega=D_{eg}E_0/\hbar$, 消去中间态得有效耦合$\hbar\Omega^2/(4\Delta)$, 二光子Rabi频率$\Omega^2/(2|\Delta|)$. 在旋转表象中$\bra fH_{\rm eff}\ket i=\hbar\Omega^2\e^{\ii\Delta\omega t}/(4\Delta)$, 单次原始$H'$直接$i\leftrightarrow f$元为零, 原题所问应指二阶有效扰动.

(b) 相干场复振幅和的模方给$U=\hbar\Omega^2[1+\cos(2kz+\Delta\omega t)]/(2\Delta)$. 常数项只移能, 光栅的$\e^{2\ii kz}/2$傅里叶分量给$\bra fU\ket i=\hbar\Omega^2\e^{\ii\Delta\omega t}/(4\Delta)$, 与(a)相同. 因而受激光力和两光子动量转移是同一过程的两种表述.
\end{solution}

\originalsection{作业十: 双光子自发辐射与暗态}
\source{Assignment 10, 第1题.}
\begin{exercise}上态$b$到下态$a$双光子发射,$\omega_1+\omega_2=W=\omega_{ba}$. (a) 从二阶吸收振幅求速率及共振条件, 验证吸收与受激发射率相等, 解释两个中间态时间顺序, 把经典场强换成模式光子数再换成模式密度, 推导Göppert--Mayer谱. 原题Gaussian单位结果为
\[A(\omega_1)=\frac{8e^4\omega_1^3\omega_2^3}{\pi\hbar^2c^6}\left\langle\left|\sum_f\left[\frac{\langle a|\mathbf r\cdot\hat e_1|f\rangle\langle f|\mathbf r\cdot\hat e_2|b\rangle}{\omega_2+\omega_{fb}}+\frac{\langle a|\mathbf r\cdot\hat e_2|f\rangle\langle f|\mathbf r\cdot\hat e_1|b\rangle}{\omega_1+\omega_{fb}}\right]\right|^2\right\rangle_{\rm avg}.\]
对$S\to S$方向平均后可写前因子$8e^4/(3\pi\hbar^2c^6)$乘$|\sum_fz_{af}z_{fb}[(\omega_1+\omega_{fb})^{-1}+(\omega_2+\omega_{fb})^{-1}]|^2$. (b) 估氢$2S$寿命: 只保留$2P$中间态, 两个$z$矩阵元取$a_0$, $2S,2P$视简并. 求$\tau^{-1}=\frac12\int_0^WA(\omega)\dd\omega$, 用秒给值并画完整谱. 真值0.122 s仅供比较.\end{exercise}
\begin{solution}(a) 时间振幅的$[\e^{\ii\delta t}-1]/\delta$模方除以$t$趋$2\pi\delta(\delta)$, 共振为$W=\omega_1+\omega_2$. 速率为$\pi E_1^2E_2^2e^4/(8\hbar^4)$乘题给中间态和的平方及共振delta. Hermiticity和能量关系$\omega_1-\omega_{fa}=-(\omega_2+\omega_{fb})$把吸收式换成发射式, 模方相等. 两项是先作用模式1、后作用模式2及反顺序; 必须相干相加后平方, 不能把两概率直接相加.

Gaussian量子场的相干振幅满足$E_i^2=8\pi\hbar\omega_i n_i/V$. 自发发射把$n_1n_2$因子替成两模式空场的$(n_1+1)(n_2+1)=1$. 三维含偏振模式密度$V\omega_i^2/(\pi^2c^3)$, 再对方向与偏振平均, 与上述激发率相乘, $V$消去并给题面前因子. 用共振delta消去$\omega_2$积分. 最后$1/2$只在总对数中除去同一光子对的双计数, 不在谱定义中再重复除一次.

(b) 简化后的$A(\omega)=C W^2\omega(W-\omega)$,$C=8e^4a_0^4/(3\pi\hbar^2c^6)$, 因分母和为$W/[\omega(W-\omega)]$. 积分$\int_0^W\omega(W-\omega)\dd\omega=W^3/6$, 得$\tau^{-1}=CW^5/12=2\alpha^2a_0^4W^5/(9\pi c^4)$, $W=3E_h/(8\hbar)$. 数值约$0.305$ s, 只表示题设极粗单中间态模型, 与真值相差数倍不能叫精确计算. 谱在两端为零, $\omega=W/2$最大且对称.
\begin{center}\begin{tikzpicture}\begin{axis}[width=.7\linewidth,height=3.8cm,xlabel={$\omega_1/W$},ylabel={归一化谱形},domain=0:1,samples=80,ymin=0]
\addplot[thick,structurecolor]{4*x*(1-x)};\end{axis}\end{tikzpicture}\end{center}
\end{solution}
\source{Assignment 10, 第2题.}
\begin{exercise}$g,e,f$三态在旋波交互表象中$H/\hbar=\bigl(\begin{smallmatrix}0&\omega_1/2&0\\\omega_1/2&\Delta&\omega_2/2\\0&\omega_2/2&-\delta\end{smallmatrix}\bigr)$. (a) $\delta=0$证明暗态不混激发态且光移零, 与$\Delta$无关; (b) $\Delta=\delta=0$缓变两耦合把暗态由$g$转至$f$; (c) 给暗亮态导数耦合的绝热条件.\end{exercise}
\begin{solution}(a) $\ket D=(\omega_2\ket g-\omega_1\ket f)/Q$, $Q=\sqrt{\omega_1^2+\omega_2^2}$, 代入Hamiltonian的$e$分量相消. (b) 令$\tan\theta=\omega_1/\omega_2$, 则$D=\cos\theta\ket g-\sin\theta\ket f$. 先开$\omega_2$, 再开$\omega_1$, 最后关闭$\omega_2$, $\theta:0\to\pi/2$. (c) 亮态为$[\sin\theta\ket g+\cos\theta\ket f\pm\ket e]/\sqrt2$, 能量$\pm\hbar Q/2$, 导数投影为$|\dot\theta|/\sqrt2$. 故$|\dot\theta|/\sqrt2\ll Q/2$, 等价于$|\omega_2\dot\omega_1-\omega_1\dot\omega_2|\ll Q^3/\sqrt2$. 端点若两个场同时为零, 能隙也为零, 不能仅称“缓慢”就保证绝热.
\end{solution}
